\documentclass[12pt]{article}
\usepackage[utf8]{inputenc}
\usepackage[T1]{fontenc}
\usepackage{hyperref}
\title{Wik-Keyangan Language Resources for K-12 Curriculum Development}
\author{Compiled by Claude Sonnet 5.5 (claude-sonnet-5-5) for Andr\'as Kornai}
\date{}
\begin{document}
\maketitle

\section{Introduction}
Wik-Keyangan (ISO 639-3: wif; Glottocode: wikk1239) is listed as an unattested, extinct Wik language of Cape York Peninsula, Queensland. AIATSIS lists it as Y173. Catalogues give variant names such as Wik-Keyenganh and rate it critically endangered. I found no speaker figures and no linguistic description.

\section{Primary Educational Resources}
No primer, textbook or school programme for Wik-Keyangan turned up in my searches.

\section{Linguistic and Reference Materials}
\begin{itemize}
\item AIATSIS record for Wik Ngathan (Y54), which lists Wik-Keyangan (Y173) in the same family, citing Sutton (1993): \url{https://collection.aiatsis.gov.au/node/1510875}
\item Wikipedia article on Wik-Me'nh, which mentions wif as unattested: \url{https://en.wikipedia.org/wiki/Wik-Me%27nh_language}
\end{itemize}

\section{Religious and Community Resources}
None specific to Wik-Keyangan found.

\section{Digital and Audio Resources}
\begin{itemize}
\item UNESCO World Atlas of Languages entry: \url{https://ar.wal.unesco.org/en/languages/wik-keyangan}
\item Endangered Languages Project entry: \url{https://endangeredlanguages.com/lang/3348?hl=en}
\end{itemize}
These are catalogue entries only. I did not open them, so I cannot say whether they hold any language data.

\section{Curriculum Development Context}
There is essentially no documentation. A school unit could only place the name within the Wik language group and work with Wik community members on what is known locally.
\end{document}